% Generated draft: requires mathematical and source review.
\documentclass[11pt]{article}
\usepackage[margin=25mm]{geometry}
\usepackage{fontspec}
\setmainfont{Latin Modern Roman}
\usepackage{amsmath,amssymb,mathtools,unicode-math}
\setmathfont{Latin Modern Math}
\usepackage{enumitem,needspace,xurl,hyperref}
\setlength{\parindent}{0pt}
\setlength{\parskip}{5pt}
\setlength{\emergencystretch}{3em}
\newcommand{\N}{\mathbb{N}}
\newcommand{\Z}{\mathbb{Z}}
\newcommand{\Q}{\mathbb{Q}}
\newcommand{\R}{\mathbb{R}}
\newcommand{\C}{\mathbb{C}}
\newcommand{\F}{\mathbb{F}}
\newcommand{\A}{\mathbb{A}}
\newcommand{\PP}{\mathbb{P}}
\newcommand{\E}{\mathbb{E}}
\newcommand{\eps}{\varepsilon}
\newcommand{\dd}{\,\mathrm{d}}
\begin{document}
\section*{UnsolvedMath ID: 800825}
\textbf{Source problem code:} catalogue entry
\subsection{Definitions and mathematical statement}
The supplied source poses the following research problem on non-orientable surfaces; it does not specify a single surface or a complete dynamical model.

For terminology, let \(M\) be a smooth non-orientable surface equipped with a positive area density \(\mu\). Interpreting area preservation as preservation of this density gives
\[
G=\operatorname{Diff}_{\mu}(M)=\{\varphi\in\operatorname{Diff}(M):\varphi^*\mu=\mu\}.
\]
This density interpretation is an explicit formulation choice: the source does not specify the area structure, and a non-orientable surface does not admit a globally defined nowhere-vanishing ordinary area two-form. Compactness, boundary conditions, additional geometric data, and the choice of the full group or a subgroup remain unspecified.

Once a Lie algebra \(\mathfrak g\) of the chosen group, an appropriate dual space \(\mathcal P\), and a class \(\mathcal A\) of admissible functionals have been fixed, a Casimir is a functional \(C\in\mathcal A\) satisfying
\[
\{C,F\}(m)=0
\qquad\text{for every }F\in\mathcal A\text{ and every }m\in\mathcal P,
\]
where \(\{\cdot,\cdot\}\) is the chosen Lie--Poisson bracket. The source does not specify these spaces or regularity classes.

\begin{enumerate}
\item Study the first non-integrable cases involving a small number of point vortices on non-orientable surfaces. Determine the relevant small-vortex configurations and describe their non-integrable dynamics after specifying the surface, vortex model, vortex strengths, Hamiltonian, and any symmetry reduction. If integrability is intended in the Liouville sense, on a symplectic phase space of dimension \(2n\) this means the existence of \(n\) first integrals, including the Hamiltonian, that are functionally independent on the stipulated regular set and pairwise Poisson commute. The source does not specify this meaning of integrability, the regularity or global domain of the integrals, or how the word ``first'' is to be ordered.
\item Describe the full set of Casimirs for the area-preserving diffeomorphism group of a non-orientable surface, with the dual space and admissible functional class made explicit.
\item Describe finite-dimensional approximations for this group or its associated hydrodynamical structures, analogous to those available in orientable cases. The approximation target, preserved structures, and required sense of convergence must be specified; the supplied record does not identify them or the orientable constructions to be generalized.
\end{enumerate}

These qualifications remain unresolved in the supplied record. The definitions above clarify terminology conditionally and do not select a particular surface, vortex system, dual space, or approximation scheme.
\subsection{Short English statement}
Study the earliest non-integrable systems of a few point vortices on non-orientable surfaces, classify the Casimirs of their area-preserving diffeomorphism groups, and describe finite-dimensional approximations analogous to orientable constructions. The source leaves the surfaces, dynamical models, integrability criterion, functional setting, and approximation requirements unspecified.
\subsection{Sources}
\begin{itemize}
\item Local UnsolvedMath registry: permanent ID \(800825\), source problem code \texttt{catalogue entry}, titled \emph{Geometric Hydrodynamics in Open Problems --- Problem 24}. The supplied registry attributes the problem to Khesin, Misiolek, and Modin, with proposed year 2022, and labels it \texttt{open}; this label is supplied curation, not an independently verified mathematical or literature assessment.
\item Khesin, Misiolek, and Modin, \emph{Geometric Hydrodynamics in Open Problems} (2022), Problem 24. Supplied source URL: \url{https://arxiv.org/abs/2205.01143}. The background records access on 2026-10-16 and extraction by \texttt{arxiv\_latex}; these are supplied provenance details, not a claim of source access or verification in preparing this draft.
\item The supplied registry records creation and update at \texttt{2026-10-16T14:00:00Z} and upstream revision \texttt{d371ee917ac9647fbb06cd6fc529deec6f3dda82}. No separate literature-verification assessment is supplied.
\end{itemize}
\end{document}
